\section{Introduction}
\label{sec:intro}

\paragraph{Why attention decoding matters.}
About $1.5$ billion people live with some hearing loss, and the number is
expected to reach $2.5$ billion by $2050$ \citep{who2021hearing}. The hardest
situation for them is the one \citet{cherry1953} called the cocktail party
problem: several people talking at once. A hearing aid can make sound louder, but
it cannot tell which voice matters, so it amplifies all of them, including the
ones the listener is trying to ignore. \emph{Auditory attention decoding} (AAD)
addresses this \citep{osullivan2015,geirnaert2021}. The brain's response follows
the speech a listener attends to \citep{ding2012,lalor2009}, so a decoder can
read that response from the electroencephalogram (EEG), find the attended
talker, and steer a hearing device toward them \citep{pan2024neuroheed}.

\paragraph{Why several streams matter.}
Such a device will not see the brain alone. Glasses and earbuds already carry
motion sensors, and eye trackers and cameras are being added to them. These
signals are useful. In everyday listening, people turn their head and eyes toward
the talker they follow \citep{hendrikse2019,best2023gaze,kidd2005}, and a hearing
aid steered by an eye tracker has already been built \citep{kidd2013visually}. A
realistic decoder should therefore combine these streams with the EEG, and doing
so raises accuracy (Section~\ref{sec:res-hinge}).

\paragraph{Why the EEG must still carry the decision.}
Gaze and head motion show where a listener is \emph{facing}, not whom they are
\emph{listening to}, and the two often differ. A driver listens to a passenger
while watching the road. At a dinner table, people listen to one person while
looking at another. Some listeners cannot move their eyes or head freely at all.
A decoder that leans on the behavioural streams does well when facing and
listening agree, and fails in exactly these cases. Those are the cases where a
brain-based device is needed: when facing and listening agree, an eye tracker is
enough. Whatever the EEG adds, it adds through the brain signal. A multimodal
decoder should therefore use the other streams as helpers, but it must not let
them replace the EEG. Accuracy cannot tell us whether this holds. A multimodal
model is reported with one number, and that number does not say which input
produced it. A decoder that silently decides from gaze and head motion and a
decoder that truly reads attention from the EEG can therefore report the same
accuracy, even though only the second keeps working when facing and listening
disagree. The risk is real: \citet{rotaru2024} found that decoders believed to
read spatial attention from EEG owed most of their accuracy to eye movements
captured in the recording.

The usual formulation is a matching task. The model sees a window of the
recording and $K$ candidate speech signals, one of which is the attended talker,
and it must pick that one; accuracy is compared with chance, $1/K$
\citep{accou2021,monesi2020,puffay2023}. We study it on MAESTRO
\citep{maestro2026}, a public dataset in which sixteen listeners each attend one
of four talkers speaking at once while EEG, eye tracking, head motion and scene
video are recorded together (Section~\ref{sec:setup}). We show that this task
can be solved without using the brain signal at all, in three ways
(Section~\ref{sec:failures}).
\textbf{(i) The audio gives the answer away.} Scaling each candidate to the same
loudness does not make the candidates look alike: statistics such as skew and
sparsity do not change with scale. On this dataset they differ between attended
and ignored talkers, and a simple classifier that sees only the audio reaches
\R{probe/raw/K4} four-way accuracy, against a chance of \R{probe/raw/K4/chance}.
\textbf{(ii) The standard scoring function gives the model a path to the leak.}
Problem (i) says the answer is \emph{present} in the audio; this one says the
standard architecture can \emph{reach} it without the brain. If the brain
encoder outputs the same vector at every time step --- the collapse an encoder
falls into when its input stops helping --- the usual score reduces to a
classifier on the audio alone, and training often finds this solution because it
is easier than the intended one.
\textbf{(iii) Easy streams crowd out the brain.} When several streams are
available, training leans on the easiest one
\citep{wang2020gradblend,huang2022competition,wu2022greedy,peng2022ogm}. Here
the easy streams measure where the listener is \emph{looking}, not what they are
\emph{listening to}, so accuracy rises while the evidence for the claim shrinks.

One test catches all three problems. We shuffle a recorded stream across the
test windows, so each window gets another window's recording but keeps its own
candidates and its own label, and we measure how much accuracy drops. We call the
drop the \textbf{contribution} of that stream. A decoder that does not use a
stream loses nothing when it is shuffled. This paper turns that test from an
occasional check into the centre of the method.

MAESTRO was released with a benchmark, and the benchmark already reports a
permutation test: it shuffles the EEG, gaze, head motion and video
\emph{together}, and its contributions are positive. That shows the listener's signals are used. It cannot show \emph{which}
signal is used, and problem (iii) sits exactly there, because a joint test gives
its largest value when a strong behavioural stream has pushed the brain aside.
\merit{} shuffles each stream \emph{separately} and splits the total exactly
across streams. It also decodes better than the benchmark:
\R{v2hinge2/H4_nohinge_nodrop/within/acc2} four-way accuracy within listeners and
\R{v2hinge2loso/H4_nohinge_nodrop/loso/acc2} on unseen listeners, against the
benchmark's $0.698$ and $0.674$ (Section~\ref{sec:res-full}).

\paragraph{Contributions.}
\textbf{(1) A per-stream measure of credit} (Section~\ref{sec:shapley}). We
compute each stream's Shapley value over all $2^M$ subsets of the $M$ recorded
streams ($M\!=\!5$ here), keeping the streams in a subset intact and shuffling
the rest.
The values add up exactly to the total contribution, and they are cheap enough
to report for every run.
\textbf{(2) A decoder in which the shortcuts are blocked by design}
(Section~\ref{sec:method}). A time-centred correlation score gives every
candidate the same score when the brain encoder is constant, so that shortcut
lands exactly at chance. The classifier heads that read gaze and head motion
receive no audio, so they cannot use audio cues. The candidate audio is tested before training.
\textbf{(3) Choosing checkpoints by contribution, not accuracy}, which trades a
little accuracy for a larger share of credit on the brain.
\textbf{(4) Findings.} Adding behavioural streams cuts the brain's share of the
credit almost in half; at two-second windows the EEG earns almost none of it.
On two other public datasets, the same tests find a label imbalance in a widely
used benchmark and show that high accuracy on both cannot be confirmed as brain
activity. We also tested adding the permutation test to the training loss
(Section~\ref{sec:hinges}); it does not increase the brain's share, and we report
this in full.
